% =====================================================================
% Relativistic Thermal Instability with Magnetic Field: §41--44
% Agent 17: Perturbation equations and stationary convection
% =====================================================================
%
% This file provides the relativistic generalization of Chapter IV,
% §41--44 of Chandrasekhar's \emph{Hydrodynamic and Hydromagnetic
% Stability}.  The classical Boussinesq MHD perturbation equations
% are replaced by a covariant Israel--Stewart formulation that
% respects causality; the Chandrasekhar number $Q$ acquires
% Alfv\'en-speed corrections; the variational principle and
% exact solutions (two free / two rigid boundaries) are re-derived;
% and every perturbation mode is checked against the causal bound.
%
% Requires: rel_preamble.tex (loaded by rel_main.tex)
% =====================================================================

\section{The problem of relativistic thermal instability in the
  presence of a magnetic field: general considerations}
\label{sec:rel-41}

The classical arguments of \S\,41 carry over almost verbatim to the
relativistic setting, but two new features appear.

\begin{enumerate}
\item \textbf{Alfv\'en speed bound.}  In the Newtonian theory the
  Alfv\'en speed $V_{A}=H\sqrt{\mu/4\pi\rho}$ can be arbitrarily
  large.  Relativistically, one must use
  \begin{equation}\label{eq:rel-vA}
    \vA^{2}
    = \frac{\bsq/4\pi}{\enthalpy/c^{2}+\bsq/(4\pi c^{2})}\,c^{2}
    = \frac{\bsq\,c^{2}}{4\pi\enthalpy+\bsq},
  \end{equation}
  where $\enthalpy=\edensity+p$ is the relativistic enthalpy density
  and $\bsq=b^{\mu}b_{\mu}$.  This satisfies $\vA<c$ for all
  finite field strengths, so the frozen-in theorem is replaced by
  the statement that perturbations propagate as relativistic Alfv\'en
  waves whose group velocity cannot exceed~$c$.

\item \textbf{Causal heat transport.}  The buoyancy--dissipation
  balance that sets the critical temperature gradient now involves
  the Israel--Stewart heat-flux relaxation time $\tauq$.  Because
  heat pulses travel at a finite speed
  $v_{q}=\sqrt{\thermcond T/(\enthalpy\,\tauq)}$, the possibility
  of overstability (oscillatory onset) acquires a causal upper
  frequency bound, in contrast to the Newtonian case where
  overstability modes can, in principle, have unbounded frequency.
\end{enumerate}

\noindent
The net effect of these two constraints is that an imposed vertical
magnetic field still inhibits convection---indeed more strongly than
in the Newtonian limit, because the Alfv\'en speed saturates---but
overstable modes are restricted to a finite band in frequency space.

%----------------------------------------------------------------------
\section{Relativistic perturbation equations}
\label{sec:rel-42}

\subsection*{(a) \emph{The background state}}

Consider an infinite horizontal layer of relativistic,
electrically-conducting fluid of depth~$d$, with a uniform vertical
magnetic field $B^{\mu}=(0,0,0,B)$ in the fluid rest frame.
The unperturbed state is static ($u^{\mu}=(\gamma c,\mathbf{0})$
with $\gamma=1$), with energy density~$\edensity_{0}$,
pressure~$p_{0}$, enthalpy density $\enthalpy_{0}=\edensity_{0}+p_{0}$,
and a steady adverse temperature gradient
\begin{equation}\label{eq:rel-beta}
  \beta = -\frac{dT_{0}}{dz}>0.
\end{equation}

The total energy--momentum tensor of the magnetized fluid is
\begin{equation}\label{eq:rel-Tmunu-MHD}
  T^{\mu\nu}_{\text{MHD}}
  = \biggl(\frac{\enthalpy}{c^{2}}+\frac{\bsq}{4\pi c^{2}}\biggr)
    u^{\mu}u^{\nu}
  + \biggl(p+\frac{\bsq}{8\pi}\biggr)g^{\mu\nu}
  - \frac{1}{4\pi}\,b^{\mu}b^{\nu},
\end{equation}
where $b^{\mu}$ is the magnetic-field four-vector satisfying
$b^{\mu}u_{\mu}=0$.

\subsection*{(b) \emph{Linearized relativistic MHD with dissipation}}

Perturb every quantity: $u^{\mu}\to u^{\mu}_{0}+\delta u^{\mu}$,
$T\to T_{0}+\Theta$, $b^{\mu}\to b^{\mu}_{0}+\delta b^{\mu}$,
$p\to p_{0}+\delta p$, etc.  In the rest frame of the background
fluid, denote the spatial perturbation velocity
$\delta u^{i}=(u,v,w)$ and the magnetic-field perturbation
$\delta b^{i}=(h_{x},h_{y},h_{z})$.  Including both viscous and
resistive dissipation through the Israel--Stewart formalism
(see~\S\,33 of the relativistic conventions), the linearized
equations are:

\medskip
\noindent\textbf{Momentum equation} (relativistic Boussinesq analogue):
\begin{equation}\label{eq:rel-87}
  \frac{\enthalpy_{0}}{c^{2}}\,\frac{\partial\delta u^{i}}{\partial t}
  = -\frac{\partial}{\partial x^{i}}(\delta\varpi_{\text{rel}})
    + \frac{\enthalpy_{0}}{c^{2}}\,g\,\alpha_{\text{rel}}\,\Theta\,\lambda^{i}
    + \shearvisc\,\nabla^{2}\delta u^{i}
    + \frac{B}{4\pi}\,\frac{\partial h^{i}}{\partial z},
\end{equation}
where $\lambda^{i}=(0,0,1)$ is the vertical unit vector,
$\alpha_{\text{rel}}$ is the relativistic thermal expansion
coefficient (cf.\ the framework chapter), and the generalized
pressure perturbation is
\begin{equation}\label{eq:rel-88}
  \delta\varpi_{\text{rel}}
  = \frac{\delta p}{\enthalpy_{0}/c^{2}}
    + \frac{B\,h_{z}}{4\pi\,\enthalpy_{0}/c^{2}}\,.
\end{equation}
The key difference from equation~(87) of the classical text is the
replacement $\rho_{0}\to\enthalpy_{0}/c^{2}$ (relativistic inertia).

\medskip
\noindent\textbf{Induction equation} (with resistive diffusion):
\begin{equation}\label{eq:rel-89}
  \frac{\partial h^{i}}{\partial t}
  = B\,\frac{\partial\,\delta u^{i}}{\partial z}
    + \eta\,\nabla^{2}h^{i},
\end{equation}
where $\eta$ is the magnetic diffusivity.  This retains its
classical form because the ideal-MHD condition
$F^{\mu\nu}u_{\nu}=0$ and Ohm's law are covariant statements.

\medskip
\noindent\textbf{Solenoidal constraint:}
\begin{equation}\label{eq:rel-90}
  \frac{\partial h^{i}}{\partial x^{i}} = 0.
\end{equation}

\medskip
\noindent\textbf{Israel--Stewart heat equation:}
\begin{equation}\label{eq:rel-91}
  \tauq\,\frac{\partial^{2}\Theta}{\partial t^{2}}
  + \frac{\partial\Theta}{\partial t}
  = \kappa_{T}\,\nabla^{2}\Theta + \beta\,w
  + \tauq\,\beta\,\frac{\partial w}{\partial t},
\end{equation}
where $\kappa_{T}$ is the thermal diffusivity.
In the Navier--Stokes limit $\tauq\to 0$ this reduces to the
classical equation~(\ref{eq:4-91}).  The $\tauq$-terms convert the
parabolic heat equation into a \emph{hyperbolic} telegrapher-type
equation, ensuring finite propagation speed for thermal signals.

\medskip
\noindent\textbf{Continuity (incompressibility in relativistic
  Boussinesq):}
\begin{equation}\label{eq:rel-92}
  \frac{\partial\,\delta u^{i}}{\partial x^{i}} = 0.
\end{equation}

\subsection*{(c) \emph{Elimination of the pressure perturbation and
  vertical-component equations}}

Taking the curl of~(\ref{eq:rel-87}) and then the curl again
(exactly as in the classical derivation, \S\,42), we eliminate
$\delta\varpi_{\text{rel}}$ and obtain the equation for the vertical
velocity~$w$:
\begin{equation}\label{eq:rel-105}
  \frac{\enthalpy_{0}}{c^{2}}\,\frac{\partial}{\partial t}\nabla^{2}w
  = \frac{\enthalpy_{0}}{c^{2}}\,g\,\alpha_{\text{rel}}
    \biggl(\frac{\partial^{2}}{\partial x^{2}}
           +\frac{\partial^{2}}{\partial y^{2}}\biggr)\Theta
  + \shearvisc\,\nabla^{4}w
  + \frac{B}{4\pi}\,\frac{\partial}{\partial z}\,\nabla^{2}h_{z}.
\end{equation}
For the vertical magnetic-field perturbation,
\begin{equation}\label{eq:rel-102}
  \frac{\partial h_{z}}{\partial t}
  = B\,\frac{\partial w}{\partial z}+\eta\,\nabla^{2}h_{z}.
\end{equation}
The vorticity ($\zeta$) and current ($\xi$) equations retain
their classical forms~(\ref{eq:4-103})--(\ref{eq:4-104}) with
$\rho_{0}\to\enthalpy_{0}/c^{2}$.

\subsection*{(d) \emph{Normal-mode analysis}}

Seek solutions proportional to
$\exp\{i(k_{x}x+k_{y}y)+pt\}$ with $k^{2}=k_{x}^{2}+k_{y}^{2}$.
Define the non-dimensional variables (in units of~$d$):
\begin{equation}\label{eq:rel-118}
  a = kd,\qquad
  \sigma = \frac{pd^{2}}{\nu_{\text{rel}}},\qquad
  p_{1} = \frac{\nu_{\text{rel}}}{\kappa_{T}},\qquad
  p_{2} = \frac{\nu_{\text{rel}}}{\eta},
\end{equation}
where $\nu_{\text{rel}}=\shearvisc\,c^{2}/\enthalpy_{0}$ is the
relativistic kinematic viscosity.
The normal-mode equations become:
\begin{align}
  \bigl(\tauq\,\sigma^{2}\nu_{\text{rel}}/d^{2}
        + p_{1}\sigma+a^{2}-D^{2}\bigr)\,\Theta
  &= \frac{\beta d^{2}}{\kappa_{T}}
     \bigl(1+\tauq\,\sigma\,\nu_{\text{rel}}/d^{2}\bigr)\,W,
  \label{eq:rel-119}\\[4pt]
  (D^{2}-a^{2}-p_{2}\sigma)\,K
  &= -\frac{Bd}{\eta}\,DW,
  \label{eq:rel-120}\\[4pt]
  (D^{2}-a^{2})(D^{2}-a^{2}-\sigma)\,W
  + \Qrel\,D(D^{2}-a^{2})\,K
  &= \frac{g\,\alpha_{\text{rel}}\,d^{2}}{\nu_{\text{rel}}}\,
     a^{2}\,\Theta,
  \label{eq:rel-121}
\end{align}
where $D=d/dz$ (with $z$ measured in units of~$d$) and the
\textbf{relativistic Chandrasekhar number} is
\begin{equation}\label{eq:rel-Qrel-def}
  \boxed{
  \Qrel
  = \frac{B^{2}d^{2}}{4\pi\,(\enthalpy_{0}/c^{2})\,\nu_{\text{rel}}\,\eta}
  = \frac{B^{2}d^{2}}{4\pi\,\rdensity\,\nu\,\eta}
    \cdot\frac{1}{1+(\edensity_{0}+p_{0}-\rdensity c^{2})/(\rdensity c^{2})}\,.
  }
\end{equation}
In the Newtonian limit $(\edensity_{0}+p_{0})\to\rdensity c^{2}$,
so $\Qrel\to Q$ of the classical theory (equation~(\ref{eq:4-134})).

\begin{relcorrection}
The relativistic Chandrasekhar number is related to the classical
one by
\[
  \Qrel
  = Q\,\frac{\rdensity c^{2}}{\enthalpy_{0}}
  = Q\,\biggl(1-\frac{\vA^{2}}{c^{2}}\biggr)
    \relcorr{\vA^{2}/c^{2}},
\]
where $\vA$ is the Alfv\'en speed~(\ref{eq:rel-vA}).
For weakly relativistic plasmas ($\vA\ll c$) the correction is
$\order{\vA^{2}/c^{2}}$.
\end{relcorrection}

The equations for the $z$-components of vorticity ($Z$) and
current density ($X$) are formally independent and identical in
structure to equations~(\ref{eq:4-122})--(\ref{eq:4-123}) with
$\Qrel$ replacing~$Q$.

\subsection*{(e) \emph{Boundary conditions for relativistic MHD}}

The boundary conditions on the hydrodynamic variables ($W$, $\Theta$,
$Z$) are exactly those of the non-relativistic theory (Chapter~II,
\S\,9\,(a)) since they express kinematic constraints that are
frame-independent.

For the magnetic perturbations, the two cases of \S\,42\,(a) carry
over:

\medskip
\noindent\textbf{Case~A} (non-conducting boundary):
\begin{equation}\label{eq:rel-107}
  X = 0,\qquad
  \mathbf{h}\text{ continuous with external vacuum field}.
\end{equation}

\noindent\textbf{Case~B} (perfectly conducting boundary):
\begin{equation}\label{eq:rel-110}
  K = 0,\qquad DX = 0,\qquad
  \nabla^{2}h_{z}=0\;\text{on the boundary}.
\end{equation}

These conditions remain the same relativistically, because they
derive from the covariant junction conditions on $F^{\mu\nu}$ across
a material interface.

%----------------------------------------------------------------------
\section{Stationary convection --- relativistic variational
  principle}
\label{sec:rel-43}

\subsection*{(a) \emph{Marginal-state equations}}

Setting $\sigma=0$ (stationary marginal state) in
equations~(\ref{eq:rel-119})--(\ref{eq:rel-121}), the
Israel--Stewart terms drop out (the $\tauq$-terms are proportional
to~$\sigma$), and we recover the \emph{same} structure as the
classical equations~(\ref{eq:4-128})--(\ref{eq:4-132}), but with
$\Qrel$ replacing~$Q$ and the relativistic Rayleigh number
\begin{equation}\label{eq:rel-Rarel-def}
  \Rarel
  = \frac{g\,\alpha_{\text{rel}}\,\beta\,d^{4}}
         {\nu_{\text{rel}}\,\kappa_{T}}
\end{equation}
replacing~$R$.

Eliminating $K$ from equation~(\ref{eq:rel-121}) via
equation~(\ref{eq:rel-120}) (with $\sigma=0$), we obtain the
single equation for~$W$:
\begin{equation}\label{eq:rel-133}
  (D^{2}-a^{2})^{2}W - \Qrel\,D^{2}W
  = \frac{g\,\alpha_{\text{rel}}\,d^{2}}{\nu_{\text{rel}}}\,
    a^{2}\,\Theta,
\end{equation}
and after applying $(D^{2}-a^{2})$ and using the heat equation,
\begin{equation}\label{eq:rel-135}
  (D^{2}-a^{2})\bigl[(D^{2}-a^{2})^{2}-\Qrel\,D^{2}\bigr]\,W
  = -\Rarel\,a^{2}\,W.
\end{equation}
This is formally identical to equation~(\ref{eq:4-135}) with
$R\to\Rarel$ and $Q\to\Qrel$.

\begin{relcorrection}
The formal identity of the marginal-state equation with the
classical one is a consequence of the Israel--Stewart terms
vanishing at $\sigma=0$.  All relativistic corrections to the
\emph{stationary} onset therefore enter exclusively through
the modified dimensionless parameters $\Rarel$ and $\Qrel$.
\end{relcorrection}

As in the classical theory, $X=Z=0$ for stationary convection.

\subsection*{(b) \emph{A variational principle}}

The variational formulation of \S\,43\,(a) of the classical text
carries over without change.  Defining
\begin{equation}\label{eq:rel-141}
  F = (D^{2}-a^{2})^{2}W - \Qrel\,D^{2}W,
\end{equation}
the Rayleigh number is given by the Rayleigh quotient
\begin{equation}\label{eq:rel-147}
  \Rarel
  = \frac{\displaystyle
      \int_{0}^{1}\bigl[(DF)^{2}+a^{2}F^{2}\bigr]\,dz}
    {a^{2}\displaystyle
      \int_{0}^{1}\bigl\{[(D^{2}-a^{2})W]^{2}
                         +\Qrel\,(DW)^{2}\bigr\}\,dz}\,.
\end{equation}
The critical $\Rarel$ for stationary convection is the
\emph{absolute minimum} of this ratio over admissible trial
functions.

\subsection*{(c) \emph{Thermodynamic significance: relativistic
  entropy production}}

In the Newtonian theory (\S\,43\,(b) of the classical text), the
variational principle expresses the balance
\begin{equation}
  \epsilon_{v}+\epsilon_{\eta} = \epsilon_{\text{buoy}},
\end{equation}
where $\epsilon_{v}$ is viscous dissipation, $\epsilon_{\eta}$ is
Joule heating, and $\epsilon_{\text{buoy}}$ is the rate of energy
liberation by buoyancy.

Relativistically, each dissipative channel contributes to the
covariant entropy production rate
$\nabla_{\mu}s^{\mu}\ge 0$ (second law).  One has
\begin{equation}\label{eq:rel-entropy-prod}
  T\,\nabla_{\mu}s^{\mu}
  = \frac{\pi^{\mu\nu}\pi_{\mu\nu}}{2\,\shearvisc}
    + \frac{\bulkP^{2}}{\bulkvisc}
    + \frac{q^{\mu}q_{\mu}}{\thermcond\,T}
    + \frac{J^{\mu}J_{\mu}}{\sigma_{\text{cond}}},
\end{equation}
where $\sigma_{\text{cond}}$ is the electrical conductivity.
In the linearized, marginal state ($\bulkP=0$, incompressible):
\begin{itemize}
  \item The shear-stress term gives the relativistic viscous
    dissipation $\epsilon_{v}^{\text{rel}}$, identical in form
    to~(\ref{eq:4-148}) with $\rho\nu\to\shearvisc$.
  \item The current term gives the relativistic Joule dissipation
    $\epsilon_{\eta}^{\text{rel}}$, identical to~(\ref{eq:4-158})
    with $Q\to\Qrel$.
  \item The heat-flux term is second-order in perturbations
    and drops out at linear order.
\end{itemize}
The balance $\epsilon_{v}^{\text{rel}}+\epsilon_{\eta}^{\text{rel}}
=\epsilon_{\text{buoy}}^{\text{rel}}$ again yields
expression~(\ref{eq:rel-147}).

\begin{relcorrection}
The variational principle retains its thermodynamic meaning
relativistically: it expresses the minimum entropy-production
condition at marginal stability, with both viscous and Ohmic
dissipation channels included in the covariant entropy
current~(\ref{eq:rel-entropy-prod}).
\end{relcorrection}

%----------------------------------------------------------------------
\section{Solutions for stationary convection}
\label{sec:rel-44}

Because the marginal-state equation~(\ref{eq:rel-135}) is formally
identical to the classical equation~(\ref{eq:4-135}) with the
substitutions $R\to\Rarel$, $Q\to\Qrel$, all classical solutions
carry over with these replacements.

\subsection*{(a) \emph{Two free boundaries}}

The eigenfunction is $W=A\sin\pi z$ and the characteristic equation
is
\begin{equation}\label{eq:rel-163}
  \Rarel
  = \frac{\pi^{2}+a^{2}}{a^{2}}
    \bigl[(\pi^{2}+a^{2})^{2}+\pi^{2}\Qrel\bigr].
\end{equation}
Letting $x=a^{2}/\pi^{2}$,
\begin{equation}\label{eq:rel-165}
  \Rarel
  = \pi^{4}\,\frac{1+x}{x}
    \biggl[(1+x)^{2}+\frac{\Qrel}{\pi^{2}}\biggr].
\end{equation}
The minimum over~$x$ occurs when
\begin{equation}\label{eq:rel-166}
  2x^{3}+3x^{2} = 1+\Qrel/\pi^{2},
\end{equation}
exactly the same cubic as in the classical theory and the
rotational problem (equation~III\,(131)), with
$\Qrel/\pi^{2}$ replacing $T/\pi^{4}$.

\subsubsection*{The $\pi^{2}Q$ law --- does it persist
  relativistically?}

For $\Qrel/\pi^{2}\gg 1$, the asymptotic root of~(\ref{eq:rel-166})
is $x_{\min}\to(\Qrel/2\pi^{2})^{1/3}$, giving
\begin{equation}\label{eq:rel-168}
  \boxed{
  \Rarel^{(c)} \;\to\; \pi^{2}\,\Qrel
  \qquad (\Qrel\to\infty),
  }
\end{equation}
and
\begin{equation}\label{eq:rel-168b}
  a_{\min} \to \bigl(\tfrac{1}{2}\pi^{4}\Qrel\bigr)^{1/6}.
\end{equation}
\emph{The $\pi^{2}Q$ law persists in the relativistic theory.}

Translating back to physical quantities:
\begin{equation}\label{eq:rel-169}
  g\,\alpha_{\text{rel}}\,\beta_{c}
  \;\to\;
  \pi^{2}\,\frac{B^{2}\,\sigma_{\text{cond}}}{4\pi\,
    (\enthalpy_{0}/c^{2})}\,\kappa_{T}\,d^{-4}
  \qquad (B\to\infty).
\end{equation}
Comparing with the classical result~(\ref{eq:4-169}):
\begin{equation}\label{eq:rel-169-compare}
  \frac{\beta_{c}^{\text{rel}}}{\beta_{c}^{\text{class}}}
  = \frac{\rdensity}{\enthalpy_{0}/c^{2}}
  = 1 - \frac{\vA^{2}}{c^{2}} + \order{\vA^{4}/c^{4}}.
\end{equation}
The relativistic critical gradient is slightly \emph{lower} (the
field is slightly \emph{more} effective at suppressing convection)
because the relativistic inertia $\enthalpy_{0}/c^{2}>\rdensity$
makes it harder for buoyancy to overcome the magnetic tension.

\begin{relcorrection}
The classical result that the critical gradient becomes independent
of viscosity in the strong-field limit remains valid
relativistically.  The only change is $\rho\to\enthalpy/c^{2}$,
giving a correction of order $\vA^{2}/c^{2}$ to the critical
temperature gradient.
\end{relcorrection}

\subsection*{(b) \emph{Two rigid boundaries}}

The variational method of \S\,44\,(b) applies without modification.
Expanding
\[
  F = \sum_{m}A_{m}\cos\bigl[(2m+1)\pi z\bigr],
\]
the secular determinant is
\begin{equation}\label{eq:rel-187}
  \biggl|\frac{1}{2}
    \biggl(\frac{c_{2n+1}}{\Rarel\,a^{2}}-\gamma_{2n+1}\biggr)
    \delta_{nm}-(n|m)\biggr| = 0,
\end{equation}
where $c_{2n+1}$, $\gamma_{2n+1}$, and $(n|m)$ are defined
as in equations~(\ref{eq:4-184}), (\ref{eq:4-176}), and
(\ref{eq:4-189}) with $Q\to\Qrel$.  The $q_{j}$ are roots of
$(q^{2}-a^{2})^{2}-\Qrel\,q^{2}=0$, i.e.\
\begin{equation}\label{eq:rel-178}
  q_{1,2}
  = \tfrac{1}{2}\bigl[\sqrt{\Qrel+4a^{2}}\pm\sqrt{\Qrel}\,\bigr].
\end{equation}

Numerical solutions of~(\ref{eq:rel-187}) at successive
approximation orders reproduce the classical Table~XV values when
$\Qrel=Q$.  For relativistic corrections, the mapping
$\Qrel = Q\,\rdensity c^{2}/\enthalpy_{0}$ shifts every entry
in the $(Q,R_{c})$ table.

The asymptotic relation $\Rarel^{(c)}\to\pi^{2}\Qrel$ holds also
for rigid boundaries, confirming the universality of the $\pi^{2}Q$
law.

\subsection*{(c) \emph{One rigid and one free boundary}}

As in \S\,44\,(c), the solution for one rigid and one free boundary
is obtained from the rigid--rigid solution by considering odd
eigenfunctions.  An odd solution for $W$ at cell depth~$d$ provides
a mixed-boundary solution at depth~$\tfrac{1}{2}d$, with
\[
  \Rarel^{\text{(mixed)}}
  = \tfrac{1}{16}\,\Rarel^{\text{(rigid--rigid)}},\qquad
  \Qrel^{\text{(mixed)}}
  = \tfrac{1}{4}\,\Qrel^{\text{(rigid--rigid)}}.
\]
This reduction is purely geometric and is unaffected by
relativistic corrections.

\subsection*{(d) \emph{Cell patterns with magnetic inhibition}}

Since $Z=0$ (no vertical vorticity) for stationary convection, the
horizontal velocity components are determined from $W(z)$ alone,
exactly as in the simple B\'enard problem.  The cell patterns
(rolls, hexagons) are identical to those of Chapter~II, \S\,16.

The \emph{quantitative} change is that the critical wave number
$a_{c}(\Qrel)$ increases with $\Qrel$, causing cells to become
narrower and more elongated.  The asymptotic scaling
\begin{equation}\label{eq:rel-acrit}
  a_{c} \sim \bigl(\tfrac{1}{2}\pi^{4}\Qrel\bigr)^{1/6}
  \qquad (\Qrel\gg 1)
\end{equation}
shows that the cell width $\ell\sim d/a_{c}$ shrinks as
$\Qrel^{-1/6}$.  Relativistically, $\Qrel<Q$ (cf.\
equation~(\ref{eq:rel-Qrel-def})), so for the same magnetic
field strength, cells are slightly \emph{wider} than the
classical prediction---a consequence of the increased inertia.

%----------------------------------------------------------------------
\subsection*{(e) \emph{Causality verification for all perturbation
  modes}}

\begin{causalitycheck}
We verify that all perturbation mode phase speeds respect the
causal bound $v_{\text{phase}}\le c$.

\begin{enumerate}
\item \textbf{Stationary convection ($\sigma=0$).}
  The marginal state is time-independent; $v_{\text{phase}}=0<c$.
  \checkmark

\item \textbf{Alfv\'en--vorticity modes.}
  The coupled vorticity--current equations
  (\ref{eq:4-122})--(\ref{eq:4-123}) with $\Qrel$ have
  characteristic speeds
  \[
    v_{\text{A,pert}}^{2}
    = \frac{\Qrel\,\nu_{\text{rel}}\,\eta}{d^{2}}
    = \frac{B^{2}}{4\pi\,\enthalpy_{0}/c^{2}}
    = \vA^{2} < c^{2}. \quad\checkmark
  \]

\item \textbf{Thermal modes.}
  The Israel--Stewart heat equation~(\ref{eq:rel-91}) is a
  telegrapher equation with characteristic speed
  \[
    v_{q}^{2}
    = \frac{\kappa_{T}}{\tauq}
    \le c^{2}
    \quad\text{(guaranteed by the Israel--Stewart constraint
    $\tauq\ge\kappa_{T}/c^{2}$).}  \quad\checkmark
  \]

\item \textbf{Magnetosonic modes.}
  The fast magnetosonic speed in the relativistic theory is
  \[
    \vf^{2}
    = \cs^{2}+\vA^{2}-\frac{\cs^{2}\vA^{2}}{c^{2}}
    < c^{2}. \quad\checkmark
  \]
  This is an algebraic identity; no additional constraint is
  needed.

\item \textbf{Overstable modes} (to be analyzed in \S\,46).
  For complex $\sigma=\sigma_{r}+i\sigma_{i}$, the group velocity
  of the perturbation is bounded by $\max(\vA,\vf,v_{q})<c$.  The
  Israel--Stewart relaxation times $\tauq$, $\taupi$ ensure that
  the dispersion relation has no superluminal roots.
  \checkmark
\end{enumerate}

\noindent\textbf{Conclusion:} All perturbation modes in the
relativistic thermal-instability problem with a vertical magnetic
field propagate at speeds strictly less than~$c$.  The theory is
manifestly causal.
\end{causalitycheck}

%----------------------------------------------------------------------
\subsection*{Summary of relativistic modifications (\S\S\,41--44)}

\begin{enumerate}
  \item The perturbation equations acquire the Israel--Stewart
    heat-relaxation term ($\tauq$-dependent), converting the heat
    equation from parabolic to hyperbolic.
  \item All inertial terms replace $\rho_{0}$ with
    $\enthalpy_{0}/c^{2}$.
  \item The Chandrasekhar number becomes
    $\Qrel = Q\,\rdensity c^{2}/\enthalpy_{0}$, incorporating
    Alfv\'en-speed corrections.
  \item For stationary convection ($\sigma=0$), the
    Israel--Stewart terms vanish, and the eigenvalue problem is
    formally identical to the classical one with
    $(R,Q)\to(\Rarel,\Qrel)$.
  \item The $\pi^{2}Q$ law $\Rarel^{(c)}\to\pi^{2}\Qrel$
    persists for strong fields.
  \item The critical temperature gradient is reduced by a factor
    $\rdensity c^{2}/\enthalpy_{0}$ relative to the classical
    value.
  \item All modes satisfy the causal bound
    $v_{\text{phase}}<c$.
\end{enumerate}
